\documentclass[10pt,a4paper]{amsart}
\usepackage[margin=19mm]{geometry}
\usepackage{amsmath,amssymb,mathtools}
\usepackage[scr=rsfs,cal=euler]{mathalfa}
\usepackage{xfrac}
\usepackage[unicode,hidelinks,bookmarks=false]{hyperref}
\newcommand{\ie}{i.e.\ }
\newcommand{\cf}{cf.\ }
\newcommand{\resp}{respectively\ }
\newcommand{\Subsection}{\subsection}
\providecommand{\nobreakdash}{\nobreak\hbox{-}}
\setlength{\emergencystretch}{2em}
\multlinegap0pt
\usepackage[shortlabels]{enumitem}
\usepackage{amssymb}
\usepackage{dsfont}
\usepackage{mathtools}
\newtheorem{assumption}{Assumption}
\newtheorem{lemma}{Lemma}
\newcommand{\E}{\mathbb{E}}
\newcommand{\Lemm}{Lemma}
\newcommand{\bA}{\mathbf{A}}
\newcommand{\bI}{\mathbf{I}}
\newcommand{\bK}{\mathbf{K}}
\newcommand{\bP}{\mathbf{P}}
\newcommand{\cD}{\mathcal{D}}
\newcommand{\cJ}{\mathcal{J}}
\newcommand{\cK}{\mathcal{K}}
\newcommand{\ind}[1]{\mathds{1}_{#1}}
\newcommand{\mbf}{\mathbf{f}}
\newcommand{\slo}{\underline{\sigma}}
\title{\LARGE \bf Training Dynamics and Induced Preconditioning \\ of ReLU Policy Gradient for Scalar LQR}
\author{Jhojan A. Rodriguez-Gil, C\'esar A. Uribe}
\date{Source: arXiv:2604.22138v2 (CC BY 4.0)}
\begin{document}
\maketitle

\noindent\emph{Source and licence.} \LARGE \bf Training Dynamics and Induced Preconditioning \\ of ReLU Policy Gradient for Scalar LQR. Version arXiv:2604.22138v2 (CC BY 4.0), distributed by arXiv under the Creative Commons Attribution 4.0 International licence, http://creativecommons.org/licenses/by/4.0/, as recorded in the arXiv licence metadata for this exact version and shown on the abstract page https://arxiv.org/abs/2604.22138v2. Full licence notice: \texttt{licenses/arxiv-2604.22138-CC-BY-4.0.txt}.\par\smallskip

\noindent\textit{Editorial scope.} Counted unit: the article's lemma lemma\_J\_geometry (source0.tex lines 347--397). The article's complete original proof of it is delivered with it (source0.tex lines 399--454, one \texttt{proof} environment of 56 lines). No mathematics is written, repaired or re-derived: the statement and the proof are the article's own lines, in the article's own notation, and no retained line is substituted or re-flowed.  Retained uncounted as the source-local context the proof needs: lines 163--165; lines 193--196; lines 198--200; lines 202--209; lines 284--305; lines 337--344.  Nothing was omitted from any retained span, so the package carries no omission record.  No retained line contains a citation, so the wrapper carries no bibliography and no bibliography entry.  No retained line contains a figure environment, an image-inclusion command or drawing code, so no image is delivered and the package declares no asset dependency.  Editorial formatting only, in addition: an AMS \texttt{amsart} preamble that restates the article's own declarations and macros used by the retained lines, namely the environments \texttt{assumption}, \texttt{lemma} on separate counters, together with the standard packages those lines need.  The retained environments print in this document as Assumption 1, Lemma 1 in order of appearance, whereas the article prints the same items with the numbering of its own full document.  The complete native source of the article as distributed in the arXiv e-print for this exact version is bundled unchanged as \texttt{sources/199073/source0.tex}.
\begin{equation} \label{eq_dynamics}
    x_{t+1} = a x_t + b u_t, \quad t=0,1,2,\dots,
\end{equation}

\begin{equation}\label{eq_riccati}
    P^{\star}=q+\frac{\gamma a^{2}rP^{\star}}{r+\gamma b^{2}P^{\star}},
    \qquad K^{\star}=-\frac{\gamma abP^{\star}}{r+\gamma b^{2}P^{\star}}.
\end{equation}

\begin{assumption} \label{assumption_controllability}
    System \eqref{eq_dynamics} is controllable, i.e., $b \neq 0$; we have access to a stabilizing controller $K_{0}$, i.e., $| a + b K_{0} | < 1$; and the optimal policy $K^{\star}$ is stabilizing, i.e., $|a + b K^{\star}| < 1$.
\end{assumption}

\begin{assumption} \label{assumption_initial_distribution}
    The moments
    $\sigma_{1}^{2} \triangleq \E[x_{0}^{2}\ind{\{x_{0} \ge 0\}}]$ and
    $\sigma_{2}^{2} \triangleq \E[x_{0}^{2}\ind{\{x_{0} < 0\}}]$
    are finite and strictly positive. We write
    $\slo^{2} {\triangleq} \min\{\sigma_{1}^{2}, \sigma_{2}^{2}\} {>} 0$ and
    $\sigma_{\cD}^{2} {\triangleq} \sigma_{1}^{2} + \sigma_{2}^{2} {=} \E[x_{0}^{2}]$.
\end{assumption}

\begin{lemma} \label{lemma_state_value_function}
   Let Assumptions~\ref{assumption_controllability}--\ref{assumption_initial_distribution} hold and let $\delta \in (0, 1)$. If $\bK\in \cK_{\overline{\delta}}$, then, writing $a_i\triangleq a+bK_i$, $J(x_{0}) = P_{1} x_{0}^2 \ind{\{x_0 \ge 0\}} + P_{2} x_{0}^2 \ind{\{x_0 < 0\}}$ where $P_{1}, P_{2}$ are uniquely defined as the solution of $\bP(\bK) = (\bI - \gamma \bA(\bK) )^{-1} \mbf(\bK)$, with
    {\small
    \begin{equation}
        \bP(\bK) {\triangleq}
        \begin{bmatrix}
            P_{1} \\
            P_{2}
        \end{bmatrix},\;\;
        \mbf(\bK) {\triangleq}
        \begin{bmatrix}
            q + rK_{1}^2 \\
            q + rK_{2}^2
        \end{bmatrix},\;\;
        \bA(\bK) {\triangleq}
        \begin{bmatrix}
            (a_{1})_{+}^2 & (a_{1})_{-}^2 \\
            (a_{2})_{-}^2 & (a_{2})_{+}^2
        \end{bmatrix}.
    \end{equation}
    }Moreover, $\cJ(\bK) = \sigma_{1}^{2} P_{1} + \sigma_{2}^{2} P_{2}$, $P_{1}, P_{2} \ge q$, and $\max(P_{1}, P_{2}) \le P_{\max} \triangleq \varphi_{\max} / (1 - \gamma(1 - \overline{\delta})^2)$ with $K_{\max} \triangleq (|a| + 1 - \overline{\delta}) / |b|$ and $\varphi_{\max} \triangleq q + rK_{\max}^2$.
\end{lemma}

\begin{equation}\label{eq_occupancy_gradient}
    \begin{gathered}
    g_i\triangleq\partial_{K_i}\cJ(\bK)
      =2d_i(\bK)\bigl(rK_i+\gamma b\,a_iP_{n(i)}\bigr)
      \\
      =2d_i(\bK)R_i(\bK)\bigl(K_i-\widehat K_i(\bK)\bigr),
    \end{gathered}
\end{equation}

\begin{lemma}\label{lemma_J_geometry}
    Let Assumptions~\ref{assumption_controllability}--\ref{assumption_initial_distribution} hold, let $\delta\in(0,\delta_0)$, and set $\bK^\star \triangleq (K^{\star}, K^{\star})$, where $P^{\star} > 0$ is the corresponding Riccati coefficient, so that $\bK^{\star} \in \cK_{\overline{\delta}}$. Then:
 
    (i) (Smoothness) There exists $L_\bK<\infty$ such that
    \begin{equation}
        \| \nabla_{\bK} \cJ(\bK) {-} \nabla_{\bK} \cJ(\bK')\|
        {\le}
        L_\bK \|\bK {-} \bK'\|,
        \; \forall\, \bK, \bK' {\in} \cK_{\overline{\delta}}.
    \end{equation}
 
    (ii) (Quadratic growth) For all $\bK \in \cK_{\overline{\delta}}$,
    \begin{equation}
        \cJ(\bK)-\cJ(\bK^\star)
        \ge
        (r+\gamma b^2P^\star)\,\slo^{2} \,\|\bK - \bK^\star\|^2.
    \end{equation}
 
    (iii) (Gradient dominance) Let $R^{\star}\triangleq r+\gamma b^2P^{\star}$ and $\mathbf d^{\star}\triangleq\mathbf d(\bK^{\star})$. Then, it holds that
      \begin{equation}
        \tfrac{1}{2} \| \nabla_{\bK} \cJ(\bK)\|^2
        \ge
        \alpha_{\bK} \bigl( \cJ(\bK) - \cJ(\bK^\star) \bigr),
        \quad
        \forall \bK \in \cK_{\overline{\delta}},
    \end{equation}
    where
    \begin{equation}\label{eq_explicit_pl}
        \alpha_{\bK}\triangleq
        \frac{2R^{\star}}{\max_{i=1,2}d_i^{\star}/(\sigma_i^2)^2}>0.
    \end{equation}
  
 
    (iv) (Uniform gradient bound) With $d_0 \triangleq 1 - \gamma(1 - \overline{\delta})^{2}$,

    \vspace{-0.2cm}
    \begin{equation}
        \|\nabla_{\bK} \cJ(\bK)\| {\le} G_{\max}
        {\triangleq}
        \frac{2\sqrt{2}\,\sigma_{\cD}^2 \bigl(rK_{\max}+\gamma |b|(1-\overline\delta)P_{\max}\bigr)}{ d_0}.
    \end{equation}
    \vspace{-0.2cm}
 
    Consequently, with $C_{\mathrm{grad}} \triangleq {L_\bK^2}/{((r{+}\gamma b^2P^\star)\slo^{2})}$,
    \begin{equation} \label{eq_cgrad}
        \| \nabla_{\bK} \cJ(\bK)\|^2
        \le
        C_{\mathrm{grad}} \bigl( \cJ(\bK) - \cJ(\bK^\star) \bigr),
        \quad \forall \bK \in \cK_{\overline{\delta}}.
    \end{equation}
\end{lemma}

\begin{proof}
\emph{(i) Smoothness.} By \Lemm~\ref{lemma_state_value_function}, $\cJ(\bK)=\sigma_{1}^{2}P_{1}(\bK) + \sigma_{2}^{2}P_{2}(\bK)$, and since $(\cdot)_+^2$ and $(\cdot)_-^2$ have Lipschitz derivatives and $\|\bA(\bK)\|_\infty\le (1 - \overline{\delta})^2$ on $\cK_{\overline\delta}$, the map $\bK \mapsto \bP(\bK)$ is $\mathcal{C}^{1,1}$ on $\cK_{\overline{\delta}}$. Hence $\nabla_{\bK} \cJ$ is Lipschitz on $\cK_{\overline{\delta}}$.
Writing $\partial_i=\partial_{K_i}$, an explicit, conservative choice is obtained by setting
\vspace{-0.1cm}
\[
\begin{aligned}
\kappa&=1-\overline\delta, \ T={2(rK_{\max}+\gamma|b|\kappa P_{\max})}/{d_0},\\
H&=\frac{2r+2\gamma b^2P_{\max}+4\gamma|b|\kappa T}{d_0}, \
L_{\bK}{=}\max\{\alpha_{\bK},\,2\sigma_{\cD}^2H\}.
\end{aligned}
\]
Differentiation of the value equations gives the following:
$\|\partial_i\bP\|_\infty\le T$ and
$\|\partial_i\partial_j\bP\|_\infty\le H$ on each sign cell.
The matching first derivatives across cell boundaries give the same
Lipschitz bound on the entire safe box.

\emph{(ii) Quadratic growth.} Completing the square in~\eqref{eq_riccati} gives
\[
(\bI-\gamma\bA)(\bP-P^\star\mathbf1)
 =R^\star[(K_1-K^\star)^2,(K_2-K^\star)^2]^\top.
\]
Multiplying by $[\sigma_1^2,\sigma_2^2](\bI-\gamma\bA)^{-1}$ yields
\vspace{-0.1cm}
\begin{equation}\label{eq_exact_gap}
\cJ(\bK)-\cJ(\bK^\star)
 =R^\star\sum_{i=1}^2d_i(\bK)(K_i-K^\star)^2.
\end{equation}
\vspace{-0.3cm}

\noindent Since $d_i(\bK)\ge\sigma_i^2\ge\slo^2$, this proves (ii).

\emph{(iii) Gradient dominance.} Let $g=\nabla_{\bK}\cJ(\bK)$ and define $Q(x,u) \triangleq qx^{2}+ru^{2}+\gamma J(ax+bu)$. For a nonzero $x$ on half-line $i$, define $\phi_i(z)\triangleq Q(x,zx)/x^2$, which does not depend on the magnitude of $x$; since $P_1, P_2\ge P^{\star}$, its piecewise second derivatives are at least $2R^{\star}$ and its first derivatives match at the breakpoint, so it is $2R^{\star}$-strongly convex, and~\eqref{eq_occupancy_gradient} reads $g_i=d_i\phi_i'(K_i)$. Telescoping $J_{\bK}-J_{\bK^{\star}}$ along the optimal trajectory, with $ \Delta\triangleq\cJ(\bK)-\cJ(\bK^{\star})$ gives
\begin{equation}
    \begin{aligned}
    \Delta&=\sum_{i=1}^2 d_i^{\star}\bigl[\phi_i(K_i)-\phi_i(K^{\star})\bigr] \le\sum_{i=1}^2\frac{d_i^{\star}}{4R^{\star}d_i^2}\,g_i^2
       \le\frac{\|g\|^2}{2\alpha_{\bK}}.
    \end{aligned}
\end{equation}

Strong convexity gives $\phi_i(K_i) -\phi_i(K^{\star}) \le \phi_i'(K_i)^2 / (4R^{\star})$, and $d_i\ge\sigma_i^2$ gives the last bound. The telescoping remainder vanishes because the optimal state contracts and $J_{\bK}(x)\le P_{\max}x^2$.  We choose the smoothness constant $L_{\bK}$ large enough ($L_{\bK}\ge\alpha_{\bK}$); increasing a valid Lipschitz constant preserves (i).

\emph{(iv) Gradient bound.} Drop the argument $\bK$ and write $\partial_{i}$ for $\partial_{K_{i}}$. On $\cK_{\overline{\delta}}$ the four ingredients are
\begin{equation} \label{eq_grad_ingredients}
    \begin{gathered}
    \|(\bI{-}\gamma\bA)^{-1}\|_{\infty} \le d_{0}^{-1},
    \qquad
    \|\bP\|_{\infty} \le P_{\max},
    \\
    \|\partial_{i}\mbf\|_{\infty} \le 2rK_{\max},
    \qquad
    \|\partial_{i}\bA\|_{\infty} \le 2|b|(1{-}\overline{\delta}),
    \end{gathered}
\end{equation}
the second from \Lemm~\ref{lemma_state_value_function} and the rest from $|K_{i}| \le K_{\max}$. Differentiating $\bP = (\bI{-}\gamma\bA)^{-1}\mbf$ gives $\partial_{i}\bP = (\bI{-}\gamma\bA)^{-1}\bigl(\partial_{i}\mbf + \gamma(\partial_{i}\bA)\bP\bigr)$, so \eqref{eq_grad_ingredients} yields $\|\partial_{i}\bP\|_{\infty} \le (2/d_{0})\bigl(rK_{\max} + \gamma|b|(1{-}\overline{\delta})P_{\max}\bigr)$. Since $\cJ = \sigma_{1}^{2}P_{1} + \sigma_{2}^{2}P_{2}$, we get $|\partial_{i}\cJ| \le \sigma_{\cD}^{2}\|\partial_{i}\bP\|_{\infty}$, and the Euclidean norm over the two components gives $G_{\max}$. Finally, $\nabla_{\bK}\cJ(\bK^{\star}) {=} 0$ and $\bK^{\star}{\in} \cK_{\overline{\delta}}$, so (i) and (ii) give~\eqref{eq_cgrad}.
\end{proof}

\end{document}
